\chapter{Introduction}
\label{ch:introduction}

The High-speed Unmanned General-purpose Observatory (HUGO) is developed as a reusable test platform for high-speed aeroelastic research. Its purpose is to bridge the gap between wind tunnel testing and full scale flight testing by providing researchers and the industry a modular and comparatively low cost aircraft on which new wing concepts, control systems, and aerospace components can be tested. The main design challenge is therefore not only to produce an unmanned aerial vehicle capable of fully substituting the role of wind tunnels, but to create a flexible flying laboratory that remains safe, maintainable, and economically viable while offering variable experimental configurations.

The Midterm Report \footnote{\href{https://github.com/Teo-Cokonaj/DSE-28/blob/main/docs/Projects/28_HUGO_Midterm_Report_final.pdf}{HUGO Midterm Report} (accessed on 11.06.2026)} concluded the configuration trade-off and HUG-CFG-303 was selected as the design configuration. This concept consists of a high-wing aircraft with static modular wing and a canard port, a T-tail, fuselage-mounted propulsion, and a retractable landing gear. It was selected because it provides the required experimental versatility while avoiding the additional structural and reliability penalties introduced by other, more complicated configurations considered. The selected aircraft layout therefore forms the baseline for the detailed design work presented in this Final Report.

A central requirement is that HUGO shall operate as a testbed for different high aspect ratio wings. 
A Planform Envelope is defined, describing the range of wing geometries, such as span, sweep angle, taper, and empennage configurations that HUGO is intended to accommodate. The aircraft structure, internal and external layout and the control system are then developed around this envelope. In addition, a standard planform is defined for cases where the customer does not provide a dedicated wing. This standard configuration is used as the reference case for performance, structural, flutter, and stability analyses.

The mission requirements drive the design in several directions. HUGO must remain below a maximum take-off mass of 50 kg, carry at least 5 kg and 5 L of payload, operate on Sustainable Aviation Fuel (SAF), and complete a nominal mission of at least 30 minutes. During this mission, it must be capable of reaching Mach 0.75 for a five minute test slot, and then fly at a lower speed regime for loitering, return and landing. The aircraft must support the collection of high-quality measurement data during the tests.

The report starts with the general aspects assessing the viability of the project itself. First, the system requirements, market analysis and certification research are performed to ensure the design serves a real mission and can be legally operated. The selected operating base is Den Helder Airport, which provides access to controlled airspace and a test area over the North Sea. Technical risks are assessed and mitigated, and the certification strategy is developed around EASA requirements, including the need for a Type Certificate and Design Organisation Approval to support repeated configuration changes.

The technical design then defines the aircraft configuration and system characteristics in more detail. The external and internal layouts are developed, and the major subsystems are integrated inside the fuselage. Material options are assessed, followed by structural analyses of the wing, fuselage, and critical interfaces. The performance analysis verifies that the aircraft can complete the required flight profile, while the stability and control analysis evaluates whether the aircraft can remain controllable over the required mass distribution and configuration range. Reliability, Availability, Maintainability, and Safety characteristics are also defined to ensure that HUGO can support repeated flight testing without compromising safety or data quality.

Finally, the report addresses manufacturing, assembly, integration, and sustainability. The manufacturing approach is selected to match the low production volume and largely composite structure.The assembly and integration plan focuses on accessibility, modularity, and safe subsystem installation. The sustainability assessment considers environmental, social, and economic aspects over the aircraft life cycle.
